\section{Error Analysis}
The first run produced a large number of low-coverage edges primarily because the decomposition dictionary did not cover many development predicates. This is a useful failure mode: it identifies precisely which predicates require analysis before the metric can be trusted. The expansion step turned this diagnostic signal into new decomposition and lexical-cue resources.

After expansion, the remaining low-coverage cases fall into four broad classes. First, some edges have weak lexical evidence because the relation is paraphrased without one of the current cues. Second, some fallback decompositions are too generic: they allow scoring but do not yet represent domain-specific roles well. Third, some texts mention entities but leave the predicate relation implicit. Fourth, synthetic adjudication may have introduced noisy factor choices for predicates that require expert domain knowledge.

This behavior is central to the method. A low score is not only a scalar failure; it can be traced to a predicate, factor, role, and evidence type. This makes the metric suitable for dataset diagnostics, generation-model error analysis, and iterative improvement of semantic resources.
